\section{Runtime, tools, and workflow}
\label{sec:runtime}

Version 0.3.0 provides a shared renderer for native LuaLaTeX, the bundled external
backend, and a companion command-line tool. The native backend requires
\texttt{engine=lua} and does not require shell escape. The default external
backend uses \texttt{texlua} and requires shell escape; it also supports pdfLaTeX.
When a custom external executable is configured, it must accept the configured
arguments, write its output file, and return a successful exit status. Merged data
is supplied as JSON. Positional indices must be a permutation of 1, 2, 3.

\subsection{Supported Liquid subset}

Supported constructs include variable output, dotted and bracketed property
paths, whitespace control, \texttt{if}/\texttt{elsif}/\texttt{else}, comparisons,
\texttt{contains}, \texttt{and}/\texttt{or}, and array/object \texttt{for} loops.
Array lookup follows Liquid's zero-based indexing. Array loops expose
\texttt{forloop.index}, \texttt{index0}, \texttt{first}, \texttt{last}, and
\texttt{length}. Available filters include \texttt{upcase}, \texttt{downcase},
\texttt{capitalize}, \texttt{strip}, \texttt{strip\_newlines}, \texttt{append},
\texttt{prepend}, \texttt{replace}, \texttt{replace\_first}, \texttt{default},
\texttt{size}, \texttt{join}, \texttt{split}, and \texttt{date}. Replacement is
literal, not a Lua pattern. Unsupported tags/filters are errors.

This is a subset rather than a complete Liquid implementation. Includes,
inheritance, assignments, loop modifiers, and full expression compatibility are
not provided. Template and data strings become LaTeX source; escaping TeX special
characters remains the template author's responsibility. Short inline expressions
use normal TeX argument parsing: they must have balanced braces and cannot contain
an unescaped percent comment. Use the verbatim \texttt{liquid} environment for
multiline templates and Liquid tags. Neither inline interface should be used in
moving arguments; use \verb|\sectionliquid| for sections.

\subsection{Search paths and data merging}

\begin{preamblecode}
\liquidsetup{
  template-path = {templates/, content/},
  data-files = {base.yml, override.json},
  data.pub = publications.json,
  data = local.yml,
  data-sources = {other = another.yml}
}
\end{preamblecode}

The supplied template path is tried first, followed by configured directories in
order. The first match wins. Relative inputs are resolved from the compilation
working directory; absolute inputs are used directly. Search directories do not
change the interpretation of data paths.

Root objects in \texttt{data-files} are merged from left to right, recursively.
Later scalar/array values replace earlier values, including empty arrays.
Next, \texttt{data.NAME} files are loaded into corresponding named objects; names
are processed in sorted order. Finally the selected default or per-call source
is merged at the root and wins conflicts. A root array may be used alone for row
rendering, but root merge sources must be objects. Aliases are independent source
selectors. Adding \texttt{data-sources} preserves previously declared aliases and
the default source.

\subsection{Convenience commands}

\begin{bodycode}
% Template's first rendered line is the title; the rest is the body.
\sectionliquid[Short TOC title]{section.liquid}
% Each row receives item plus forloop metadata.
\begin{itemize}
  \liquidlist[publications]{cv.yml}{publication-row.liquid}
\end{itemize}
% Check syntax/data without inserting generated output.
\liquidvalidate{section.liquid}
% Log generated code without inserting it into the document.
\inputliquid*{section.liquid}
\end{bodycode}

The optional section title is used by the table of contents. It should be a valid
LaTeX moving argument; the generated title must fit on the first rendered line.
The optional list path selects an array within the data; omit it for a root
array. Wrong data structures produce an error. Row templates work inside tables,
itemize, or arbitrary surrounding environments. Lists use the native or bundled
renderer. Validation checks syntax in the supported subset and any configured
data files; it cannot validate a custom external renderer's additional dialect.

\subsection{Custom filters and diagnostics}

\begin{preamblecode}
\liquidsetup{engine=lua,diagnostics,verbose=trace,debug-dir=debug}
\liquidfilter{shout}{function(value)
  return string.upper(value) .. "!"
end}
\end{preamblecode}

Custom filters require the native backend and should be pure functions of their
arguments. The function body is Lua source in a normal TeX argument; avoid TeX
special characters or register more complex functions directly from a Lua module.
External state/upvalues are not automatically tracked: use \texttt{no-cache} for
those functions. The native renderer reports template paths, line information,
and a Lua traceback on failures; external failure output and exit status are
forwarded into the LaTeX log.

\texttt{verbose=false} or \texttt{off} disables informational logging;
\texttt{true}/\texttt{basic} reports renders and cache hits, \texttt{full} adds input
paths, and \texttt{trace} adds bridge/renderer commands. Errors remain visible at
all levels. \texttt{debug-dir} stores template, merged data, output, and log
snapshots with stable cache keys. \texttt{diagnostics} collects request counts,
cache-hit counts, and CPU time; \verb|\liquidstats| prints them. CPU time excludes
external child-process wall time.

POSIX quoting supports spaces and apostrophes. Windows supports ordinary paths
with spaces and rejects shell metacharacters rather than executing an ambiguous
command. Windows end-to-end verification and broader path support remain planned.

\subsection{CLI and multi-output pipelines}

\begin{terminalblock}
texlua scripts/liquidtex.lua render --template cv.liquid --data cv.yml
texlua scripts/liquidtex.lua validate --template cv.liquid --data cv.yml
texlua scripts/liquidtex.lua inspect --data cv.yml
texlua scripts/liquidtex.lua pipeline examples/pipeline.json
\end{terminalblock}

The CLI can render to standard output or \texttt{--output FILE}, validate syntax
and configured data without typesetting, and inspect merged data as JSON.
\texttt{--data-file FILE} is repeatable; \texttt{--data FILE} is applied last.
Use \texttt{--cache} to enable persistent CLI caching and \texttt{--no-cache} to
force rendering. Run \texttt{--help} for the full command summary.

A pipeline is a YAML/JSON object with an \texttt{outputs} array. Each output
requires \texttt{template} and \texttt{output}, and may set \texttt{data},
\texttt{data-files} (written \texttt{data\_files} in configuration), or
\texttt{namespaces}. Shared values can be supplied at the root. All paths are
relative to the configuration file's directory. Pipelines cache each output
independently, rerun only affected jobs when inputs change, and update destination
files only when their contents differ. See the tracked CV example for a minimal
complete configuration. These pipelines render text; PDF compilation remains the
responsibility of LaTeX or a build tool.
